\chapter{Hilbert Spaces}\label{ch:b3:hilbert}

A \omterm{def:b3:hilbert:inner}{Hilbert space} is a Banach space whose norm comes from an
\omterm{def:b3:hilbert:inner}{inner product} --- and that single extra structure restores, in
infinite dimension, almost all of Euclidean geometry:
orthogonal projections exist, every \omterm{def:b3:topology:continuity}{continuous} functional is an
\omterm{def:b3:hilbert:inner}{inner product} against a fixed vector (Riesz), and orthonormal
bases expand every vector in a convergent series with
Pythagorean bookkeeping (\omterm{thm:b3:hilbert:parseval}{Parseval}). The chapter's climax is a
debt honored: the trigonometric system is an orthonormal basis
of $L^2$, so \omterm{thm:b3:hilbert:parseval}{Parseval}'s identity holds for \emph{every}
\omterm{def:b3:lebesgue:l1}{square-integrable} function --- the statement Year 2 could only
prove for piecewise $\mathcal C^1$ functions. We end with
Lax--Milgram, the workhorse lemma of the variational approach
to differential equations.

Throughout, $H$ is a vector space over $K = \R$ or $\C$.

\section{Inner products; the projection theorem}

\begin{definition}\label{def:b3:hilbert:inner}
An \emph{inner product}\index{inner product} is a map $\langle
\cdot,\cdot\rangle \colon H\times H \to K$, linear in the
second variable, with $\langle y, x\rangle =
\overline{\langle x, y\rangle}$ and $\langle x, x\rangle > 0$
for $x \neq 0$. It induces the norm $\norm x = \langle x,
x\rangle^{1/2}$, the \emph{Cauchy--Schwarz inequality}
$\abs{\langle x, y\rangle} \leq \norm x\norm y$ (Year 2's
proof --- the discriminant --- is unchanged), and the
\emph{parallelogram law}
\[
\norm{x + y}^2 + \norm{x - y}^2 = 2\norm x^2 + 2\norm y^2 .
\]
A \emph{Hilbert space}\index{Hilbert space} is an
inner-product space \omterm{def:b3:complete:complete}{complete} for this norm. Examples: $\ell^2$
(\cref{pb:b3:banach:1}) and, the fundamental one, $L^2(\mu)$
with $\langle f, g\rangle = \int\bar fg\,\dd\mu$ ---
\omterm{def:b3:complete:complete}{complete} by Riesz--Fischer (\cref{thm:b3:lp:complete}); the
inner product is finite by Cauchy--Schwarz ($= $ H\"older at
$p = q = 2$).
\end{definition}

\begin{theorem}[Projection onto a closed convex set]
\label{thm:b3:hilbert:projection}
Let $C \neq \varnothing$ be a closed \emph{convex} subset of
the \omterm{def:b3:hilbert:inner}{Hilbert space} $H$ and $x \in H$. There is a unique $p_C(x)
\in C$ with\index{projection theorem}
\[
\norm{x - p_C(x)} = d(x, C),
\]
characterized by: $\operatorname{Re}\langle x - p_C(x),\ c -
p_C(x)\rangle \leq 0$ for all $c \in C$. The map $p_C$ is
$1$-Lipschitz.
\end{theorem}

\begin{proof}
Let $d = d(x, C)$ and $(c_n) \subseteq C$ with $\norm{x - c_n}
\to d$. Parallelogram on $x - c_n$ and $x - c_m$:
\[
\norm{c_n - c_m}^2 = 2\norm{x - c_n}^2 + 2\norm{x - c_m}^2 -
4\,\bigl\|x - \tfrac{c_n + c_m}2\bigr\|^2
\leq 2\norm{x{-}c_n}^2 + 2\norm{x{-}c_m}^2 - 4d^2
\]
(convexity puts the midpoint in $C$): the right side tends to
$0$, so $(c_n)$ is Cauchy, and its limit $p \in C$ (closed)
attains $d$. Uniqueness: two minimizers give, by the same
identity, $\norm{p - p'}^2 \leq 2d^2 + 2d^2 - 4d^2 = 0$.

Characterization: for $c \in C$, $t \in \intoc01$, the vector
$p + t(c - p) \in C$, so
\[
d^2 \leq \norm{x - p - t(c-p)}^2
= d^2 - 2t\operatorname{Re}\langle x - p, c - p\rangle +
t^2\norm{c-p}^2 ;
\]
divide by $t \to 0^+$: $\operatorname{Re}\langle x - p, c -
p\rangle \leq 0$. Conversely this inequality gives $\norm{x -
c}^2 = \norm{x - p}^2 - 2\operatorname{Re}\langle x - p, c -
p\rangle + \norm{p - c}^2 \geq \norm{x-p}^2$. Lipschitz: for
$x, y$ with projections $p, q$, add the two variational
inequalities (with $c = q$, resp.\ $c = p$):
$\operatorname{Re}\langle x - y - (p - q), p - q\rangle \geq
0$, so $\norm{p - q}^2 \leq \operatorname{Re}\langle x - y, p -
q\rangle \leq \norm{x - y}\norm{p - q}$.
\end{proof}

\begin{theorem}[Orthogonal decomposition]
\label{thm:b3:hilbert:decomposition}
Let $F$ be a \emph{closed subspace} of $H$. Then $p_F$ is
linear, $x - p_F(x) \perp F$ for all $x$, and\index{orthogonal
complement}
\[
H = F \oplus F^\perp,
\qquad F^\perp = \{y : \langle y, f\rangle = 0\ \forall f\in
F\},
\qquad (F^\perp)^\perp = F .
\]
For a general subspace, $(F^\perp)^\perp = \bar F$; in
particular $F$ is dense iff $F^\perp = \{0\}$.
\end{theorem}

\begin{proof}
For a subspace, the variational characterization with $c =
p_F(x) \pm f$ ($f \in F$, both signs, and $\iu f$ in the
complex case) forces $\langle x - p_F(x), f\rangle = 0$: the
residual is orthogonal to $F$. Decomposition $x = p_F(x) + (x
- p_F(x))$ with $F \cap F^\perp = \{0\}$ ($\langle y, y
\rangle = 0$); linearity of $p_F$ follows from uniqueness of
such decompositions (both sides linear in them). $(F^\perp)
^\perp \supseteq F$ always; conversely if $x \perp F^\perp$,
write $x = f + g$: $g = x - f \in F^\perp$ and $\langle g,
g\rangle = \langle x, g\rangle - \langle f, g\rangle = 0$: $x
= f \in F$. For a general subspace $F$: $F^\perp = \bar
F^{\,\perp}$ (\omterm{def:b3:topology:continuity}{continuity} of the \omterm{def:b3:hilbert:inner}{inner product}), so
$(F^\perp)^\perp = \bar F$ by the closed case; \omterm{ex:b3:lebesgue:gamma}{density} iff
$\bar F = H$ iff $F^\perp = 0$.
\end{proof}

\begin{theorem}[Riesz representation]\label{thm:b3:hilbert:riesz}
For every \omterm{def:b3:topology:continuity}{continuous} linear functional $\varphi \in H'$ there
is a unique $a \in H$ with\index{Riesz representation theorem}
\[
\varphi(x) = \langle a, x\rangle \quad (x \in H),
\qquad \norm\varphi_{H'} = \norm a .
\]
\end{theorem}

\begin{proof}
If $\varphi = 0$: $a = 0$. Otherwise $F = \ker\varphi$ is a
closed proper subspace; pick $u \in F^\perp$, $\norm u = 1$
(\cref{thm:b3:hilbert:decomposition}: $F^\perp \neq 0$ since
$F \neq H$). For any $x$, the vector $\varphi(x)u -
\varphi(u)x \in \ker\varphi$, hence $\perp u$:
\[
0 = \langle u, \varphi(x)u - \varphi(u)x\rangle
= \varphi(x) - \varphi(u)\langle u, x\rangle :
\qquad \varphi(x) = \langle
\overline{\varphi(u)}\,u,\ x\rangle .
\]
So $a = \overline{\varphi(u)}u$ works. Uniqueness: $\langle a
- a', x\rangle = 0$ for all $x$, test $x = a - a'$. Norms:
$\abs{\varphi(x)} \leq \norm a\norm x$ (Cauchy--Schwarz) with
equality at $x = a$.
\end{proof}

\begin{example}[A projection, computed to the end]
\label{ex:b3:hilbert:projectionexample}
In $H = L^2(\intcc01)$, what is the best approximation of
$f(x) = x^2$ by an affine function? The subspace $F =
\operatorname{Vect}(1, x)$ is closed (finite-dimensional),
and $p_F(f) = a + bx$ is characterized by orthogonality of
the residual to $1$ and to $x$:
\[
\int_0^1(x^2 - a - bx)\,\dd x = 0,
\qquad
\int_0^1x\,(x^2 - a - bx)\,\dd x = 0,
\]
i.e.\ $\frac13 = a + \frac b2$ and $\frac14 = \frac a2 +
\frac b3$: $a = -\frac16$, $b = 1$. So $p_F(x^2) = x -
\frac16$, and the error is
\[
d(f, F)^2 = \int_0^1\Bigl(x^2 - x + \frac16\Bigr)^2\dd x =
\frac1{180},
\qquad d(f, F) = \frac1{6\sqrt5} .
\]
Two remarks worth internalizing. First, the computation is
nothing but a $2\times2$ linear system --- the \emph{normal
equations}; for the monomial basis their matrix
$\bigl(\frac1{i+j+1}\bigr)$ is the notoriously
ill-conditioned Hilbert matrix, and orthogonalizing first
(Legendre polynomials, \cref{pb:b3:hilbert:1}) is the cure.
Second, the best \emph{uniform} approximation of $x^2$ by
affine functions is different ($x - \frac18$, by
equioscillation): each norm has its own geometry, and only
the Hilbertian one answers with a linear system.
\end{example}

\section{Orthonormal bases}

\begin{definition}\label{def:b3:hilbert:onb}
A family $(e_i)_{i\in I}$ is \emph{orthonormal} if $\langle
e_i, e_j\rangle = \delta_{ij}$, and a \emph{Hilbert
basis}\index{Hilbert basis} (orthonormal basis) if moreover
its finite linear combinations are dense in $H$ (the family is
\emph{total}). We treat the countable case $I = \N$, which by
Gram--Schmidt covers every \emph{separable} $H$
(\cref{prop:b3:hilbert:gramschmidt}).
\end{definition}

\begin{theorem}[Bessel, Parseval]\label{thm:b3:hilbert:parseval}
Let $(e_n)_{n\in\N}$ be orthonormal in $H$, and $c_n(x) =
\langle e_n, x\rangle$.
\begin{enumerate}
  \item (Bessel) $\sum_n\abs{c_n(x)}^2 \leq \norm x^2$, and
        the series $\sum_nc_n(x)e_n$ converges in $H$, with
        sum $p_F(x)$, $F = \overline{\operatorname{Vect}}(e_n)$.
  \item The following are equivalent: (i) $(e_n)$ is a \omterm{def:b3:hilbert:onb}{Hilbert
        basis}; (ii) $x = \sum_nc_n(x)e_n$ for every $x$; (iii)
        \emph{Parseval}\index{Parseval's identity}: $\norm x^2
        = \sum_n\abs{c_n(x)}^2$ for every $x$; (iv) the only
        vector orthogonal to all $e_n$ is $0$.
  \item If $(e_n)$ is a \omterm{def:b3:hilbert:onb}{Hilbert basis}, $x \mapsto (c_n(x))_n$
        is an isometric isomorphism $H \to \ell^2$
        (\emph{every} infinite-dimensional separable \omterm{def:b3:hilbert:inner}{Hilbert
        space} ``is'' $\ell^2$), and $\langle x, y\rangle =
        \sum_n\overline{c_n(x)}c_n(y)$.
\end{enumerate}
\end{theorem}

\begin{proof}
(1) For finite $N$: $x - \sum_{n\leq N}c_ne_n \perp e_k$ ($k
\leq N$), so Pythagoras gives $\norm x^2 = \sum_{n\leq
N}\abs{c_n}^2 + \norm{x - \sum_{n\leq N}c_ne_n}^2$: Bessel.
The partial sums $S_N = \sum_{n\leq N}c_ne_n$ are Cauchy:
$\norm{S_N - S_M}^2 = \sum_{M<n\leq N}\abs{c_n}^2$, tail of a
convergent series; the limit lies in $F$, and $x - \lim S_N
\perp$ each $e_k$ (\omterm{def:b3:topology:continuity}{continuity}), hence $\perp F$: by uniqueness
of the orthogonal decomposition, $\lim S_N = p_F(x)$.

(2) (i)$\Rightarrow$(ii): $F = H$, so $p_F = \mathrm{id}$.
(ii)$\Rightarrow$(iii): Pythagoras in the limit ($\norm{S_N}^2
= \sum_{n \leq N}\abs{c_n}^2 \to \norm x^2$).
(iii)$\Rightarrow$(iv): $x \perp$ all $e_n$ gives $\norm x^2 =
0$. (iv)$\Rightarrow$(i): $F^\perp = \{0\}$
(orthogonality to all $e_n$ is orthogonality to $F$), so $F$
is dense by \cref{thm:b3:hilbert:decomposition}; but $F$, a
\omterm{def:b3:topology:interior}{closure}, is already closed: $F = H$.

(3) The map is linear, isometric by (iii) (hence injective),
and surjective: given $(c_n) \in \ell^2$, the series $\sum
c_ne_n$ converges (Cauchy as in (1)) to a preimage. The \omterm{def:b3:hilbert:inner}{inner
product} formula is polarization from (iii), or a direct limit
computation.
\end{proof}

\begin{proposition}[Gram--Schmidt]\label{prop:b3:hilbert:gramschmidt}
Let $(x_n)$ be a linearly independent sequence. Setting
inductively $\tilde e_n = x_n - \sum_{k<n}\langle e_k,
x_n\rangle e_k$ and $e_n = \tilde e_n/\norm{\tilde e_n}$
produces an orthonormal $(e_n)$ with the same finite spans:
$\operatorname{Vect}(e_1, \dots, e_n) = \operatorname{Vect}
(x_1, \dots, x_n)$. Consequently every separable \omterm{def:b3:hilbert:inner}{Hilbert space}
(one with a countable dense subset) has a \omterm{def:b3:hilbert:onb}{Hilbert
basis}.\index{Gram--Schmidt orthonormalization}
\end{proposition}

\begin{proof}
Induction: $\tilde e_n \perp e_k$ ($k < n$) by construction,
and $\tilde e_n \neq 0$ by independence; the spans match at
each stage (triangular change of basis). For a separable $H$:
from a dense sequence extract a linearly independent
subfamily with dense span (discard each vector in the span of
its predecessors — the span is unchanged), orthonormalize:
the result is total.
\end{proof}

\begin{theorem}[The trigonometric system; Parseval at last]
\label{thm:b3:hilbert:fourier}
In $L^2(\intcc{-\pi}\pi)$ with $\langle f, g\rangle =
\frac1{2\pi}\int_{-\pi}^\pi \bar fg$, the family $e_n(t) =
\eu^{\iu nt}$, $n \in \Z$, is a \omterm{def:b3:hilbert:onb}{Hilbert basis}. Consequently,
for \emph{every} $f \in L^2$ --- in particular every piecewise
\omterm{def:b3:topology:continuity}{continuous} $2\pi$-periodic $f$ --- with $c_n(f) =
\frac1{2\pi}\int_{-\pi}^{\pi}f(t)\eu^{-\iu nt}\dd t$:
\[
f = \sum_{n\in\Z}c_n(f)\,\eu^{\iu nt} \ \ \text{in } L^2,
\qquad
\frac1{2\pi}\int_{-\pi}^{\pi}\abs f^2 =
\sum_{n\in\Z}\abs{c_n(f)}^2 .
\]
This proves, in full generality, the \omterm{thm:b3:hilbert:parseval}{Parseval} identity that
Year 2 admitted.\index{Parseval's identity (general)}
\end{theorem}

\begin{proof}
Orthonormality is a direct computation (Year 2). Totality: let
$f \in L^2$ be $\perp$ all $e_n$, i.e.\ all Fourier
coefficients vanish. \omterm{def:b3:topology:continuity}{Continuous} $2\pi$-periodic functions are
dense in $L^2(\intcc{-\pi}\pi)$: indeed $\mathcal
C_c(\intoo{-\pi}\pi)$ is dense
(\cref{thm:b3:lp:density}(2)) and such functions extend
periodically and \omterm{def:b3:topology:continuity}{continuously}. Trigonometric polynomials are
$\norm\cdot_\infty$-dense among \omterm{def:b3:topology:continuity}{continuous} periodic functions
(Stone--Weierstrass, \cref{cor:b3:complete:weierstrass}(c)),
and $\norm\cdot_2 \leq \norm\cdot_\infty$: trigonometric
polynomials are dense in $L^2$. But $f \perp$ every
trigonometric polynomial, hence $f \perp$ a dense subspace: $f
\in (\text{dense})^\perp = \{0\}$
(\cref{thm:b3:hilbert:decomposition}). Criterion (iv) of
\cref{thm:b3:hilbert:parseval} concludes; (ii) and (iii)
unpack to the display (reindexing the countable $\Z$; the
double-ended series converges unconditionally --- the partial
sums over any exhausting family converge, by the $\ell^2$ tail
argument).
\end{proof}

\begin{theorem}[Lax--Milgram]\label{thm:b3:hilbert:laxmilgram}
Let $H$ be a real \omterm{def:b3:hilbert:inner}{Hilbert space} and $a \colon H\times H \to
\R$ bilinear, \emph{\omterm{def:b3:topology:continuity}{continuous}} ($\abs{a(u,v)} \leq M\norm
u\norm v$) and \emph{coercive} ($a(u, u) \geq \alpha\norm u^2$,
$\alpha > 0$). Then for every $\varphi \in H'$ there is a
unique $u \in H$ with\index{Lax--Milgram theorem}
\[
a(u, v) = \varphi(v) \qquad \text{for all } v \in H .
\]
\end{theorem}

\begin{proof}
For fixed $u$, $v \mapsto a(u, v)$ is \omterm{def:b3:topology:continuity}{continuous} linear: Riesz
gives a unique $Au \in H$ with $a(u,v) = \langle Au,
v\rangle$; $A$ is linear with $\norm{Au} \leq M\norm u$
(uniqueness of representatives, then bound). Coercivity:
$\alpha\norm u^2 \leq a(u,u) = \langle Au, u\rangle \leq
\norm{Au}\norm u$, so $\norm{Au} \geq \alpha\norm u$: $A$ is
injective with closed range (a Cauchy image sequence $Au_n$
forces $u_n$ Cauchy). The range is dense: $w \perp
\operatorname{im}A$ gives $0 = \langle Aw, w\rangle \geq
\alpha\norm w^2$. Closed and dense: $A$ is bijective. Given
$\varphi$, let $f$ represent it (Riesz) and $u = A^{-1}f$:
$a(u, v) = \langle f, v\rangle = \varphi(v)$, uniquely
($a(u - u', \cdot) = 0$ and coercivity).
\end{proof}

\begin{remark}\label{rem:b3:hilbert:variational}
When $a$ is symmetric, Lax--Milgram's solution is the unique
minimizer of the \emph{energy} $J(v) = \frac12a(v,v) -
\varphi(v)$ (\cref{exo:b3:hilbert:9}): existence of solutions
to variational problems in one stroke. Applied to suitable
function spaces (the Sobolev spaces of a later course), this
solves boundary value problems for differential equations ---
the modern entry point to partial differential equations.
\end{remark}

\section{Exercises}

\begin{exercise}[$\star$]\label{exo:b3:hilbert:1}
(a) Prove the polarization identities (real: $4\langle x,
y\rangle = \norm{x+y}^2 - \norm{x-y}^2$; complex: the
four-term version).
(b) Show that $\norm\cdot_1$ on $L^1(\intcc01)$ and
$\norm\cdot_\infty$ on $\mathcal C(\intcc01)$ violate the
parallelogram law: these norms come from no \omterm{def:b3:hilbert:inner}{inner product}.
\end{exercise}

\begin{exercise}[$\star$]\label{exo:b3:hilbert:2}
In $H = L^2(\intcc01)$ (real):
(a) compute the projection of $f$ onto the subspace of
constant functions, and interpret;
(b) compute the projection onto $\{g : g = 0 \text{ a.e.\ on }
\intcc0{1/2}\}$;
(c) compute $d\bigl(x \mapsto x,\ \operatorname{Vect}(\mathbf
1)\bigr)$.
\end{exercise}

\begin{exercise}[$\star\star$]\label{exo:b3:hilbert:3}
(a) Show that for a subspace $F$: $F$ dense $\iff$ $F^\perp =
\{0\}$, and give an example in $\ell^2$ of a \emph{proper}
dense subspace (so $F^\perp = 0$ without $F = H$: the
decomposition theorem genuinely needs $F$ closed).
(b) Show that if $x_n \to x$ and $y_n \to y$ in norm, then
$\langle x_n, y_n\rangle \to \langle x, y\rangle$, and locate
two places where the chapter used this \omterm{def:b3:topology:continuity}{continuity}.
\end{exercise}

\begin{exercise}[$\star\star$]\label{exo:b3:hilbert:4}
Apply Gram--Schmidt to $1, x, x^2$ in $L^2(\intcc{-1}1)$
(\omterm{def:b3:measure:lebesgueouter}{Lebesgue measure}): obtain the first three normalized
\emph{Legendre polynomials}, and verify they match
$\sqrt{n + \frac12}\,P_n$ for the Rodrigues polynomials $P_n$
of \cref{pb:b3:hilbert:1}.
\end{exercise}

\begin{exercise}[$\star\star$]\label{exo:b3:hilbert:5}
Apply \omterm{thm:b3:hilbert:parseval}{Parseval} (\cref{thm:b3:hilbert:fourier}) to $f(t) = t$
and $f(t) = t^2$ on $\intcc{-\pi}\pi$ --- now legitimately for
these (\omterm{def:b3:topology:continuity}{continuous}, but previously the identity needed
piecewise-$\mathcal C^1$ care at the wrap-around
discontinuity): recover
\[
\sum_{n\geq1}\frac1{n^2} = \frac{\pi^2}6,
\qquad
\sum_{n\geq1}\frac1{n^4} = \frac{\pi^4}{90} .
\]
\end{exercise}

\begin{exercise}[$\star\star$]\label{exo:b3:hilbert:6}
(a) Find $a \in L^2(\intcc01)$ with $\int_0^{1/2}f =
\langle a, f\rangle$ for all $f$; compute $\norm\varphi$ for
this functional.
(b) Show that the evaluation $f \mapsto f(\frac12)$, defined
on the subspace $\mathcal C(\intcc01) \subseteq
L^2(\intcc01)$, is \emph{not} \omterm{def:b3:topology:continuity}{continuous} for
$\norm\cdot_2$: no Riesz representative exists (evaluation is
not an $L^2$ notion).
\end{exercise}

\begin{exercise}[$\star\star\star$]\label{exo:b3:hilbert:7}
Let $H$ be separable with \omterm{def:b3:hilbert:onb}{Hilbert basis} $(e_n)$, and $(x_k)$ a
bounded sequence.
(a) Show that some subsequence converges \emph{weakly}: there
is $x$ with $\langle y, x_{k_j}\rangle \to \langle y,
x\rangle$ for every $y \in H$. \emph{(Diagonal extraction on
the coefficients $\langle e_n, x_k\rangle$; assemble $x$ via
Bessel and uniform boundedness of norms.)}
(b) Show $e_n \rightharpoonup 0$ but $\norm{e_n} = 1$: weak
limits can lose norm. Show $\norm x \leq
\liminf\norm{x_{k_j}}$ in (a).
\end{exercise}

\begin{exercise}[$\star\star$]\label{exo:b3:hilbert:8}
(Adjoints) For $T \in \mathcal L(H)$, show there is a unique
$T^* \in \mathcal L(H)$ with $\langle Tx, y\rangle = \langle
x, T^*y\rangle$ (Riesz), and $\vertiii{T^*} = \vertiii T$.
Compute the adjoint of the shift $S$ on $\ell^2$, and prove
$\ker T^* = (\operatorname{im}T)^\perp$ --- deduce
$\overline{\operatorname{im}T} = (\ker T^*)^\perp$.
\end{exercise}

\begin{exercise}[$\star\star$]\label{exo:b3:hilbert:9}
Let $a$ be as in Lax--Milgram and moreover \emph{symmetric}.
Show that $u$ solves $a(u, \cdot) = \varphi$ iff $u$ minimizes
$J(v) = \frac12a(v, v) - \varphi(v)$, and that the minimum is
attained at exactly one point. \emph{(\omterm{def:b3:complete:complete}{Complete} the square:
$J(u + w) - J(u) = \frac12a(w,w) \geq
\frac\alpha2\norm w^2$.)} Application: re-derive the
projection theorem for closed subspaces from Lax--Milgram.
\end{exercise}

\begin{exercise}[$\star\star\star$]\label{exo:b3:hilbert:10}
(The Haar system) On $\intcc01$, let $h_{0} = \mathbf 1$, and
for $n = 2^j + k$ ($j \geq 0$, $0 \leq k < 2^j$):
\[
h_n = 2^{j/2}\Bigl(\mathbf 1_{[k2^{-j},\,(k +
\frac12)2^{-j})} - \mathbf 1_{[(k+\frac12)2^{-j},\,(k+1)2^{-j})}
\Bigr).
\]
Show that $(h_n)_{n\geq0}$ is orthonormal in $L^2(\intcc01)$,
and total. \emph{(Orthogonality: disjoint or nested supports;
totality: finite spans contain all dyadic step functions,
which are dense --- via \cref{thm:b3:lp:density}(1) and dyadic
approximation of intervals.)} The Haar system is the ancestor
of wavelets.
\end{exercise}

\begin{exercise}[$\star\star$]\label{exo:b3:hilbert:11}
(Orthogonal projections, characterized) Let $H$ be a \omterm{def:b3:hilbert:inner}{Hilbert
space} and $P \in \mathcal L(H)$ with $P^2 = P$, $P \neq 0$.
Show the equivalence of:
(i) $P$ is the orthogonal projection onto
$\operatorname{im}P$;
(ii) $P = P^*$ (\cref{exo:b3:hilbert:8});
(iii) $\vertiii P = 1$.
\emph{(For (iii) $\Rightarrow$ (i): if some $x \in
(\ker P)^\perp$ had $Px \neq x$, consider $x + t(Px - x)$
--- or directly: for $u \in \operatorname{im}P$ and $v \in
\ker P$, expand $\norm{P(u + tv)}^2 \leq \norm{u + tv}^2$
for all $t \in \R$ and conclude $\langle u, v\rangle = 0$.)}
Exhibit a non-orthogonal projection on $\R^2$ and compute
its norm.
\end{exercise}

\begin{exercise}[$\star\star\star$]\label{exo:b3:hilbert:12}
(Von Neumann's ergodic theorem) Let $U \in \mathcal L(H)$ be
\emph{unitary} ($U^*U = UU^* = I$), $F = \ker(U - I)$ the
fixed space, $P$ the orthogonal projection onto $F$, and
$A_n = \frac1n\sum_{k=0}^{n-1}U^k$.
(a) Show $\ker(U - I) = \ker(U^* - I)$ \emph{(from
$\norm{Ux - x}^2 = 2\norm x^2 - 2\operatorname{Re}\langle
Ux, x\rangle$ and unitarity)}, and deduce
$\overline{\operatorname{im}(U - I)} = F^\perp$.
(b) Show that $A_nx \to x$ for $x \in F$, and $A_nx \to 0$
for $x \in \operatorname{im}(U - I)$ \emph{(telescoping)},
then for $x \in \overline{\operatorname{im}(U - I)}$
(uniform bound $\vertiii{A_n} \leq 1$).
(c) Conclude: $A_nx \to Px$ for \emph{every} $x \in H$ ---
time averages converge to the projection on the invariants.
(d) Spell it out for $H = L^2(\R/\Z)$ and $Uf = f(\cdot +
\alpha)$ with $\alpha$ irrational: identify $F$ (use
Fourier series, \cref{thm:b3:hilbert:fourier}) and deduce
that $\frac1n\sum_{k<n}f(x + k\alpha) \to \int_0^1f$ in
$L^2$: the $L^2$ equidistribution of irrational rotations.
\end{exercise}

\section{Problem: orthogonal polynomials}

\begin{problem}[{Weekend problem --- Legendre, Hermite, and
Gauss quadrature}]\label{pb:b3:hilbert:1}
Let $I \subseteq \R$ be an interval and $w > 0$ a \omterm{def:b3:topology:continuity}{continuous}
\emph{weight} on the \omterm{def:b3:topology:interior}{interior} of $I$ such that $\int_I
\abs t^nw(t)\dd t < \infty$ for all $n$; work in $H = L^2(I,
w\,\dd\lambda)$ with $\langle f, g\rangle = \int_I \bar
fg\,w$. Gram--Schmidt applied to $1, t, t^2, \dots$ produces
the \emph{\omterm{pb:b3:hilbert:1}{orthogonal polynomials}} $(p_n)$ for $w$ (monic
normalization: $p_n = t^n + \cdots$).\index{orthogonal
polynomials}

\medskip
\noindent\textbf{Part I --- General theory.}
\begin{enumerate}
  \item Show that $p_n$ is orthogonal to every polynomial of
        degree $< n$, and that $(p_0, \dots, p_n)$ is a basis
        of $\R_n[t]$.
  \item (Three-term recurrence) Show there are reals $a_n,
        b_n$ with
        \[
        p_{n+1}(t) = (t - a_n)\,p_n(t) - b_n\,p_{n-1}(t),
        \qquad b_n = \frac{\norm{p_n}^2}{\norm{p_{n-1}}^2} >
        0 .
        \]
        \emph{(Expand $t\,p_n$ in the basis $(p_k)_{k \leq
        n+1}$ and kill coefficients by orthogonality, using
        $\langle tp_n, p_k\rangle = \langle p_n,
        tp_k\rangle$.)}
  \item (Roots) Show that $p_n$ has $n$ \emph{distinct} roots,
        all \omterm{def:b3:topology:interior}{interior} to $I$. \emph{(Let $t_1 < \dots < t_m$ be
        the \omterm{def:b3:topology:interior}{interior} sign changes of $p_n$; if $m < n$,
        test $p_n$ against $\prod_{i\leq m}(t - t_i)$ and
        contradict orthogonality.)}
\end{enumerate}

\medskip
\noindent\textbf{Part II --- Legendre ($I = \intcc{-1}1$, $w =
1$).} Define $P_n(t) = \frac{1}{2^nn!}\,\frac{\dd^n}{\dd
t^n}\bigl[(t^2 - 1)^n\bigr]$ (Rodrigues).
\begin{enumerate}[resume]
  \item Show $\deg P_n = n$ with leading coefficient
        $\frac{(2n)!}{2^n(n!)^2}$, and, integrating by parts
        $n$ times, that $\langle P_n, Q\rangle = 0$ for every
        polynomial $Q$ of degree $< n$: the $P_n$ are (up to
        normalization) the \omterm{pb:b3:hilbert:1}{orthogonal polynomials} for $w = 1$.
  \item Compute $\norm{P_n}_2^2 = \frac{2}{2n+1}$
        \emph{(integrate by parts $n$ times against itself and
        reduce to a Beta/Wallis integral,
        \cref{exo:b3:product:8})}.
  \item Show that the normalized Legendre polynomials form a
        \omterm{def:b3:hilbert:onb}{Hilbert basis} of $L^2(\intcc{-1}1)$
        \emph{(Weierstrass, \cref{cor:b3:complete:weierstrass},
        plus \omterm{ex:b3:lebesgue:gamma}{density} of $\mathcal C$ in $L^2$)}, and expand
        $f(t) = \abs t$ up to degree $2$: compute the best
        quadratic $L^2$-approximation of $\abs t$.
\end{enumerate}

\medskip
\noindent\textbf{Part III --- Hermite ($I = \R$, $w(t) =
\eu^{-t^2}$).} Define $H_n(t) =
(-1)^n\eu^{t^2}\frac{\dd^n}{\dd t^n}\eu^{-t^2}$.
\begin{enumerate}[resume]
  \item Show that $H_n$ is a polynomial of degree $n$ with
        leading coefficient $2^n$, that $H_{n+1} = 2tH_n -
        H_n'$, and that $\langle H_m, H_n\rangle_w =
        \delta_{mn}\,2^nn!\sqrt\pi$ \emph{(parts again)}.
  \item Show that the Hermite family is total in $L^2(\R,
        \eu^{-t^2}\dd t)$, admitting one result from
        \cref{ch:b3:fouriertransform}: if $g \in L^1(\R)$ has
        $\int g(t)\eu^{-\iu\xi t}\dd t = 0$ for all $\xi$,
        then $g = 0$ a.e. \emph{(For $f \perp$ all $H_n$,
        i.e.\ $\perp$ all polynomials: show $z \mapsto \int
        f(t)\eu^{-t^2}\eu^{-\iu zt}\dd t$ is well defined,
        expand the exponential as a series, justify the
        interchange by domination, and conclude that the
        Fourier transform of $f\eu^{-t^2}$ vanishes.)}
\end{enumerate}

\medskip
\noindent\textbf{Part IV --- Gauss quadrature.} Fix $n$, let
$t_1 < \dots < t_n$ be the roots of $p_n$ (Part I), and define
the weights $w_i = \int_I \ell_i(t)\,w(t)\dd t$ where
$\ell_i$ are the Lagrange interpolation basis polynomials at
the $t_i$.
\begin{enumerate}[resume]
  \item Show that the quadrature rule $Q(f) = \sum_iw_if(t_i)$
        is exact on all polynomials of degree $\leq n - 1$
        (interpolation), and in fact --- the miracle --- on
        all polynomials of degree $\leq 2n - 1$: write $P =
        qp_n + r$ and use orthogonality on the quotient $q$.
  \item Show that the weights are positive
        \emph{(apply the rule to $\ell_i^2$, of degree $2n -
        2$)}, and deduce from Polya's theorem
        (\cref{exo:b3:banach:9}) that Gauss quadrature
        converges: $Q_n(f) \to \int_I fw$ for every \omterm{def:b3:topology:continuity}{continuous}
        $f$ on a \omterm{def:b3:topology:compact}{compact} $I$.
  \item For $n = 2$, $I = \intcc{-1}1$, $w = 1$: compute the
        nodes $\pm\frac1{\sqrt3}$ and weights $1, 1$, and
        verify exactness on $1, t, t^2, t^3$ by hand. Compare
        with the trapezoid rule on the same two evaluation
        points.
\end{enumerate}

\medskip
\noindent\textbf{Part V --- Chebyshev: the polynomials that
oscillate best.} Now $I = \intcc{-1}1$ and $w(t) =
\frac1{\sqrt{1 - t^2}}$.
\begin{enumerate}[resume]
  \item Show that $T_n(\cos\theta) = \cos n\theta$ defines a
        polynomial $T_n$ of degree $n$ (establish $T_{n+1} =
        2t\,T_n - T_{n-1}$ from a trigonometric identity),
        with leading coefficient $2^{n-1}$ for $n \geq 1$;
        and that the substitution $t = \cos\theta$ gives
        \[
        \langle T_m, T_n\rangle_w =
        \int_0^\pi\cos m\theta\,\cos n\theta\,\dd\theta
        = 0 \ (m \neq n), \qquad
        \norm{T_0}_w^2 = \pi,\ \ \norm{T_n}_w^2 = \frac\pi2 :
        \]
        the $T_n$ are the \omterm{pb:b3:hilbert:1}{orthogonal polynomials} for this
        weight, and Chebyshev expansions \emph{are} Fourier
        cosine series in disguise.
  \item Locate explicitly the $n$ roots $t_k =
        \cos\frac{(2k-1)\pi}{2n}$ and the $n + 1$ extrema
        $s_j = \cos\frac{j\pi}n$ of $T_n$ on $\intcc{-1}1$,
        where $T_n(s_j) = (-1)^j$: the graph
        \emph{equioscillates} between $\pm1$.
  \item (Minimax) Show that among all \emph{monic}
        polynomials of degree $n$, the polynomial
        $2^{1-n}T_n$ has the smallest sup-norm on
        $\intcc{-1}1$, namely $2^{1-n}$ --- and it is the
        unique minimizer. \emph{(If a monic $P$ had
        $\sup\abs P < 2^{1-n}$, the difference $2^{1-n}T_n -
        P$, of degree $\leq n-1$, would alternate in sign at
        the $n+1$ equioscillation points.)}
  \item Application to interpolation: for nodes $t_1 < \dots
        < t_n$ in $\intcc{-1}1$, the error of Lagrange
        interpolation of a $\mathcal C^n$ function involves
        $\omega(t) = \prod_i(t - t_i)$. Show that choosing
        Chebyshev roots as nodes minimizes
        $\sup_{\intcc{-1}1}\abs\omega$, and give the
        resulting bound $\norm{f -
        L_nf}_\infty \leq \frac{\norm{f^{(n)}}_\infty}
        {2^{n-1}\,n!}$ --- compare with equally spaced nodes
        (state Runge's phenomenon as the cautionary tale).
  \item Verify $\abs{T_n'(\pm1)} = n^2$ \emph{(differentiate
        $T_n(\cos\theta) = \cos n\theta$ and take limits
        $\theta \to 0, \pi$)}: polynomials bounded by $1$ on
        $\intcc{-1}1$ can have derivative as large as $n^2$
        at the edge (Markov's inequality says no larger ---
        statement only). Where in the interval is the
        derivative bound only $O(n)$?
  \item (Chebyshev--Gauss quadrature) Show that the Gauss
        rule for the weight $w$ at the $n$ Chebyshev roots
        has \emph{equal} weights $w_i = \frac\pi n$
        \emph{(exactness on $T_0, \dots, T_{n-1}$ plus the
        trigonometric sums
        $\sum_{k=1}^n\cos\bigl(j\tfrac{(2k-1)\pi}{2n}\bigr) =
        0$ for $1 \leq j \leq n - 1$)}: the most uniform of
        all quadratures. Write it out for $n = 3$.
\end{enumerate}

\medskip
\noindent\textbf{Part VI --- Christoffel--Darboux,
interlacing, and the Jacobi matrix.} Back to a general weight;
$h_k = \norm{p_k}^2$ (monic $p_k$), $b_k = h_k/h_{k-1}$.
\begin{enumerate}[resume]
  \item (Least norm) Show that among all \emph{monic}
        polynomials of degree $n$, the orthogonal $p_n$ is
        the unique one of minimal $L^2(w)$-norm --- identify
        the minimization as an orthogonal projection onto
        $\R_{n-1}[t]$ (\cref{thm:b3:hilbert:projection} or
        the finite-dimensional projection of Year 2). The
        minimax property of question 14 is the same statement
        with $L^\infty$ in place of $L^2$: same hero, two
        norms.
  \item (Christoffel--Darboux) Prove, by induction on $n$
        using the three-term recurrence, the identity
        \[
        \sum_{k=0}^{n}\frac{p_k(x)\,p_k(y)}{h_k}
        = \frac{p_{n+1}(x)\,p_n(y) -
        p_n(x)\,p_{n+1}(y)}{h_n\,(x - y)}
        \qquad (x \neq y),
        \]
        and its confluent form ($y \to x$):
        $\sum_{k\leq n}\frac{p_k(x)^2}{h_k} =
        \frac{p_{n+1}'(x)p_n(x) - p_n'(x)p_{n+1}(x)}{h_n}$.
  \item Deduce that $p_n$ and $p_{n+1}$ have no common root,
        and that at every root $x_0$ of $p_{n+1}$:
        $p_n(x_0)\,p_{n+1}'(x_0) > 0$. Conclude the
        \emph{interlacing} of roots: between two consecutive
        roots of $p_{n+1}$ lies exactly one root of $p_n$.
  \item (Jacobi matrix) Let $J_n$ be the $n\times n$
        symmetric tridiagonal matrix with diagonal $a_0,
        \dots, a_{n-1}$ and off-diagonal entries $\sqrt{b_1},
        \dots, \sqrt{b_{n-1}}$. Show by induction that
        $\det(tI_n - J_n) = p_n(t)$, so the roots of $p_n$
        are the eigenvalues of a real symmetric matrix ---
        re-proving in one line that they are real, and (with
        the interlacing above) tying \omterm{pb:b3:hilbert:1}{orthogonal polynomials}
        to the spectral world of \cref{ch:b3:spectral}.
  \item (Synthesis) Assemble the dictionary for the three
        classical families (Legendre, Hermite, Chebyshev):
        interval, weight, defining formula, three-term
        recurrence, norm, and the natural habitat of each
        (quadrature and approximation on compacta; Gaussian
        analysis; minimax and Fourier-cosine methods). One
        sentence on what the general theory (Parts I, VI)
        gave that no individual computation could.
\end{enumerate}

\medskip
\noindent\textbf{Part VII --- The error term, and the kernel
behind the weights.} Here $I$ is \omterm{def:b3:topology:compact}{compact} and $f \in \mathcal
C^{2n}(I)$.
\begin{enumerate}[resume]
  \item (Gauss error formula) Let $Hf$ be the Hermite
        interpolant of degree $\leq 2n - 1$ matching $f$ and
        $f'$ at the nodes $t_1, \dots, t_n$ (prove its
        existence and the pointwise error
        \[
        f(t) - Hf(t) =
        \frac{f^{(2n)}(\xi_t)}{(2n)!}\;p_n(t)^2
        \]
        by the usual auxiliary-function argument). Deduce, by
        integrating this identity against $w$ and squeezing
        between the extrema of $f^{(2n)}$, that
        \[
        \int_I f\,w - Q_n(f)
        = \frac{f^{(2n)}(\xi)}{(2n)!}\,h_n
        \qquad\text{for some } \xi \in I,
        \]
        with $h_n = \norm{p_n}^2$ as in Part VI: Gauss
        quadrature errs by one $2n$-th derivative, weighted by
        the squared norm of the monic orthogonal polynomial.
  \item (The weights are Christoffel values) Using the
        reproducing kernel $K_n(x, y) = \sum_{k=0}^{n-1}
        \frac{p_k(x)p_k(y)}{h_k}$ of $\R_{n-1}[t]$ and the
        exactness of $Q_n$ up to degree $2n - 2$, prove
        \[
        w_i \;=\;
        \Bigl(\,\sum_{k=0}^{n-1}
        \frac{p_k(t_i)^2}{h_k}\Bigr)^{\!-1} :
        \]
        each weight is the value at its node of the
        \emph{Christoffel function} --- positivity of the
        weights (question 10) again, now with an exact
        formula. Verify it recovers $w_1 = w_2 = 1$ for $n =
        2$, $I = \intcc{-1}1$, $w = 1$.
  \item (Everything checks on one integral) For the Chebyshev
        weight and $n = 3$ nodes, compute both sides of
        \[
        \int_{-1}^{1}\frac{t^6}{\sqrt{1 - t^2}}\,\dd t
        = \frac{5\pi}{16},
        \qquad
        Q_3(t^6) = \frac{9\pi}{32},
        \]
        so the quadrature error is exactly $\frac{\pi}{32}$;
        then check that the error formula of question 23
        predicts precisely this value (here $f^{(6)} = 6!$ is
        constant, and $h_3 = \norm{2^{-2}T_3}_w^2 =
        \frac\pi{32}$): theory and computation agree to the
        last digit.
\end{enumerate}
\end{problem}
